%% ==========================================================================
%% Related Work patch — three 2025 references added per Reviewer 1.C3.
%% Each is positioned against U-HAP rather than appended; insertion points
%% in the existing §II are marked with %% INSERT comments.
%% ==========================================================================

%% INSERT into §II.B (Kubernetes admission-time authorization paragraph),
%% just before the closing summary sentence. Two references go here.

Recent 2025 work has continued to push the granularity of Kubernetes
authorization. Rostamipoor~\textit{et~al.}~\cite{rostamipoor2025kubekeeper}
introduce \emph{KubeKeeper}, an admission-webhook-based mechanism that
restricts pod access to secrets at sub-RBAC granularity by intercepting
\texttt{CREATE} and \texttt{UPDATE} requests at the same extension point
U-HAP exploits; their evaluation is request-time only and does not
quantify policy-update propagation. Cesarano and
Natella~\cite{cesarano2025kubefence} present \emph{KubeFence}, which
argues that RBAC alone lacks attribute-level granularity and compiles
per-workload API filters at deploy time --- a design decision that
aligns with U-HAP's compile-then-enforce philosophy, but operates at the
container-image level rather than as a unified ACL/RBAC/ABAC/Deny
authorizer.

%% INSERT into §II.C (formal/verification paragraph). One reference here.

Closer in spirit to our compile-time correctness arguments,
Sissodiya~\textit{et~al.}~\cite{sissodiya2025formal} model Kubernetes
RBAC and admission policies as first-order predicates and use SMT
solvers to detect misconfigurations before deployment. While they
target verification rather than enforcement, their formal predicate
encoding parallels U-HAP's Phase~2 graph-and-index compilation; we view
the two as complementary --- their tool can validate inputs to U-HAP's
compiler.
